\Titular%
{divulgacion}%
{Grandes colaboracións para buscar o máis pequeno}%
{Xabier Cid Vidal}%
{}%

\begin{multicols}{2}

Luns, 9 da mañá. A semana comeza coma sempre, mirando a lista de reunións que
vou ter. Non deixo de sorprenderme pola cantidade delas ás que acudín nos anos
que levo traballando na física de partículas experimental. Posiblemente
esteamos a falar dunha media dunha reunión por día, e de máis dunha hora de
duración cada unha! E estou seguro de que o meu caso non é dos máis esaxerados.
Recordo a mesma queixa de colegas da colaboración CMS, que bromeaban sobre o
nome do seu experimento. Xa non facía referencia ao seu famoso solenoide
(\textit{Compact Muon Solenoid}), senón a que eran unha \textit{Continuous
Meeting Society} (Sociedade de Reunións Continuas). E non obstante, non hai
outro xeito. Non podemos facer física de partículas experimental sen reunirnos,
é dicir, sen colaborar.

Son membro da colaboración LHCb desde hai máis de 15 anos. LHCb é un dos catro
grandes experimentos que rexistran o resultado das colisións do LHC, o Gran
Colisor de Hadróns. Situado no CERN (Xenebra, Suíza), o LHC é o maior, máis
potente e máis enerxético colisor de partículas da historia. LHCb, cuns 1800
membros, que representan a máis de 100 centros de investigación de 28 países, é
unha colaboración pequena dentro do LHC, media para física de partículas na
actualidade, e enorme desde un punto de vista histórico. As colaboracións máis
grandes do LHC, ATLAS e CMS, responsables do descubrimento do bosón de Higgs,
contan con máis de 6000 membros cada unha na actualidade. Como referencia, as
colaboracións UA1 e UA2 do CERN, que descubriron os bosóns W e Z nos anos 80 do
século pasado, contaban cada unha con menos de 150 colaboradores. Canto máis
fundamental e pequeno é o que queremos buscar, máis persoal científico fai
falta!

E uso «persoal científico» en lugar de «físicos e físicas» intencionadamente,
pois o papel doutros campos, como o da enxeñaría, é crucial en física de
partículas. Desde a enxeñaría informática, que precisamos para manexar a
inmensa cantidade de datos coa que debemos tratar (o LHC, cando funciona, xera
máis datos que Facebook ou Instagram), ata profesionais en enxeñería de
estruturas que deseñan as cavernas nas que, a 100 m baixo terra, temos que
instalar os nosos experimentos.

Os medios de comunicación tenden a vendernos a idea da ciencia como o resultado
dun esforzo individual. É verdade que moitos dos grandes avances científicos ao
longo da historia viñeron da man de xenios incribles, e tamén o é que sen o
talento individual a ciencia e a tecnoloxía non poden progresar. Pero hoxe en
día é difícil imaxinar unha disciplina na área científico-tecnolóxica na cal
non sexa necesario o traballo en equipo. Evidentemente, o nivel de colaboración
non é igual en todas as áreas, pero estou bastante seguro de que é necesario na
actualidade un certo grao en todas elas.

\begin{figure}[H]
    \centering
    \includegraphics[width=\linewidth]{./revistas/005/imaxes/ccc.jpg}
    \caption{Council Chamber do CERN. (Foto: CERN)}
\end{figure}

En relación a esta cuestión quero mencionar a famosa serie de televisión «The
Big Bang Theory». É certo que contribuíu a popularizar o noso traballo, e,
incluso, hai quen asegura que houbo un boom no número de matriculados
universitarios en física grazas a ela. Non obstante, cando a vía, non podía
evitar sentir que o meu traballo na física de partículas non estaba ben
representado. Entendo que os creadores se han de tomar licenzas artísticas,
pero... apenas hai colaboracións! E de habelas, só ocorren entre os
protagonistas, malia que un dos personaxes principais é experimental. Por non
mencionar que o CERN só aparece referenciado de forma moi esporádica. Todos
estaremos de acordo en que calquera persoa que investigue en física de
partículas na actualidade usa a palabra CERN ou LHC varias veces á semana polo
menos.

Volvamos falar da miña colaboración, LHCb. A pesar de ter un tamaño mediano,
LHCb ten unha estrutura bastante complexa, que é un bo reflexo de como
funcionan as colaboracións en física de partículas. A colaboración divídese en
grupos relacionados cos temas de física nos que se traballa e outros máis
técnicos, relacionados con aspectos como a reconstrución, identificación de
partículas ou sistema de disparo. Ademais, existen grupos relacionados cos
diferentes sub-detectores cos que conta o experimento. Cada grupo conta
habitualmente con dúas persoas coordinadoras e os grupos de física inclúen
ademais unha subestrutura interna, con varias subcoordinacións. Todo isto é
necesario para abarcar todos os temas nos que traballamos!

\begin{figure}[H]
    \centering
    \includegraphics[width=\linewidth]{./revistas/005/imaxes/lhcb.jpg}
    \caption{Foto grupal da colaboración LHCb. (Foto: CERN)}
\end{figure}

LHCb naceu inicialmente como un experimento dedicado á física do sabor. Esta
fai referencia, por exemplo, ás dinámicas fundamentais que separan materia de
antimateria. Ademais, para iso, LHCb foi deseñado especialmente para detectar
partículas que conteñen o quark b, xa que estas son idóneas para esta misión.
Porén, como sucede ás veces, o detector resultou ser útil para moitas máis
cousas que para as que fora concibido. E iso fai que o rango de física que
cobre LHCb se ensanchase enormemente, facendo máis complexa a estrutura dos
grupos de física. Sirva como exemplo que estamos a tomar datos en colisións que
involucran ións pesados de chumbo. Aínda que o LHC colide protóns durante a
maior parte do ano, unha pequena parte do tempo colídense ións de chumbo, tanto
uns contra outros como contra protóns. Ninguén cría que LHCb estivese preparado
para estas colisións. Pero non só foi así, senón que conseguimos ter resultados
moi interesantes neste ámbito e incluso temos un Grupo de Física exclusivamente
dedicado a iso, con importante participación de persoal do IGFAE.

Para entender mellor como funciona a colaboración, poderiamos establecer un
paralelismo cunha república. A maioría das colaboracións están lideradas por
unha Portavocía, similar á presidencia, que exerce como cara pública do
experimento, e ten a misión de establecer as liñas fundamentais da
colaboración, incluíndo en que gastamos os cartos ou como nos organizamos. Hai
ademais dúas vice-portavocías, que exercen de número dous. As persoas
coordinadoras dos grupos, como explicaba anteriormente, terían un papel similar
aos ministros ou ministras. Existe ademais un posto especialmente relevante a
medio camiño entre a Portavocía e os grupos de física, que é o de Coordinador
ou Coordinadora da Física. Esta persoa coordina os grupos de física e tamén
exerce como cara visible da colaboración en moitos congresos, e tamén conta
cunha vice-coordinación. Temos a sorte de contar na USC e IGFAE con María
Vieites, investigadora Ramón y Cajal e actual vice-coordinadora da física de
LHCb.

A Portavocía, a Coordinación de Física e a Coordinación dos grupos representan
o poder executivo. Pero existe tamén un poder lexislativo, formado polos
institutos de investigación que forman parte da colaboración, denominado Xunta
da Colaboración, que debe ratificar todas e cada unha das decisións. Esta Xunta
reúnese catro veces ao ano, e cada instituto ten un voto nela. Para rematar o
paralelismo, en LHCb ata temos unha Constitución. Establece o conxunto de
normas sobre como funciona a Colaboración e foi ratificada pola Xunta. Iso si,
é unha Constitución que se cambia con relativa frecuencia. Ademais, a Xunta da
Colaboración encárgase de decidir que novos institutos poden formar parte da
mesma. As colaboracións dos experimentos do LHC non pararon de medrar nos
últimos anos.

Tal e como establece a constitución de LHCb, as posicións directivas da
colaboración escóllense por votación. Para a Portavocía, utilízase o voto
directo dos institutos, sen ningún tipo de peso. Para a Coordinación de Física,
os membros da colaboración votan individualmente. E para o resto de posicións
executivas, envíanse candidaturas, que son moi tidas en conta pola Portavocía
ou a Coordinación da Física, quen toman a decisión final. En todo caso, a
última palabra sobre os cargos tena a Xunta da Colaboración. Pero, salvo en
casos moi raros, esta só ratifica o que se votou. E este, como calquera outro
proceso democrático, implica a necesidade de movementos políticos. Existen
varios factores a ter en conta para estes cargos: nacionalidade, xénero,
institución, idade... Conseguir un cargo executivo nunha colaboración é un
aspecto importante para o currículo científico, polo que normalmente hai máis
dunha candidatura para cada posto. Copiando do inglés, falamos da necesidade de
ter «sombreiros". Loxicamente, canto máis grande o sombreiro, máis grande a
responsabilidade.

Falaba ao comezo das reunións. Son abundantes e necesarias. Pero tamén son
representativas do funcionamento dunha colaboración. Normalmente as reunións
teñen unha certa periodicidade (sobre todo semanal ou bisemanal) e teñen unha
sala física e outra virtual para a súa celebración. A sala física está no CERN,
para aqueles membros da colaboración que se atopan en Xenebra cando esta
sucede. A virtual, para o resto, habitualmente a maioría, que se atopan
dispersos polo mundo e se conectan por videoconferencia. A sala virtual faise
en plataformas en liña como Zoom, deseñadas para este tipo de eventos. Cando un
debe reunirse e entre os asistentes atópanse colaboradores de todo o mundo,
atopar unha hora adecuada pode ser todo un reto! De feito, a única hora válida
cando nos xuntamos con xente de América e de Asia é a sobremesa europea,
arredor das 2 da tarde. Implica un madrugón para os americanos e traballar ata
tarde para os asiáticos, pero é a única solución. As reunións case sempre se
desenvolven utilizando presentacións con transparencias. O nome «transparencia"
non deixa de ser unha reminiscencia do pasado, loxicamente na actualidade
utilízanse ferramentas como PowerPoint ou LaTeX. E o que máis hai nas reunións
é diálogo e discusión, moita discusión!

Toda esta estrutura é intricada e complexa, pero creo que é fundamental para
que podamos facer ciencia. Non poderiamos facela sen as colaboracións, pero non
podemos pretender coordinar a centos de científicos e científicas sen un mínimo
de organización. Gústame pensar nas nosas colaboracións como unha gran familia,
e realmente o noso grao de identificación dentro da colaboración é elevado.
Cando estamos no CERN, ou en eventos onde se reúne moito persoal investigador,
xuntámonos cos membros da nosa colaboración de forma case inconsciente, e en
moitos casos antepoñendo este criterio á nacionalidade. Nas colaboracións
traballan conxuntamente persoas de países que pode que sexan rivais, e incluso
inimigos no terreo político, o cal sería inimaxinable fóra do ámbito
científico. Do mesmo xeito que cando necesitamos atopar algo pequeno na casa,
como as chaves, toda a familia se une como un equipo para logralo, as
colaboracións de física de partículas fan o mesmo para buscar novos fenómenos e
partículas nas escalas máis pequenas. Neste caso, non é unha elección, é unha
obriga. Non poderiamos facer ciencia de alto nivel sen colaborar. En tempos tan
convulsos nos que institucións, gobernos e organismos deberían buscar o acordo
ante os grandes retos, por riba das diferenzas, as colaboracións científicas
deberían servirnos de exemplo.

\end{multicols}

% AQUI PODE POÑERSE UNHA PROPAGANDA, CARTEL, ANUNCIO OU FOTO
